% Abstract, short version
Many systems chain an AI model's typed decisions into a workflow. They
then use the model's own confidence to decide which cases a human should
check. We test that practice with Jev, TypeSafe's typed decision model.
Its three primitives are a categorical \choice{}, a binary \noul{} and an
ordinal \score{}, and we chain them into a bug-triage workflow. After
exploratory work on 442 Mozilla and Linux kernel bugs, we ran
pre-registered tests on 500 new bugs, four further models and an
independent Kubernetes workflow. A failure is a disagreement between the
workflow's action and a target action built from reference labels. That
target has not yet been validated against human judgment. A frozen
estimator that reads only the workflow's decisions ranks failures better
than final-step confidence (AUROC 0.76 vs.\ 0.55 on the new bugs). This
holds for all five models, and most of the signal is how often the
chosen action was wrong before. The Kubernetes replication agreed (AUROC
0.78 vs.\ 0.63), as did three of four further models. Only one of those
three, OpenJev, was competent at every step of the workflow. These error rates
do not transfer across label mixes. A label-free correction that helped
on the bug trackers made ranking worse on Kubernetes. Better ranking also
did not reduce cost. Under the pre-registered cost matrix, with reviewers
assumed correct, no 10\% review policy beat no review. The savings sit
in wrongly accepted cases that none of the pre-registered signals could
pick out.
